\section{Preguntas y respuestas seleccionadas: discusión de temas de vanguardia}

La sesión de preguntas y respuestas de la conferencia abordó varios temas importantes; a continuación se presenta una selección de lo más relevante:

\subsection{Agotamiento de datos y datos sintéticos}

\textbf{Pregunta}: A medida que los datos de entrenamiento se consumen continuamente y hay cada vez menos datos disponibles, ¿hacia dónde se dirige el futuro?

Puntos clave de la respuesta del profesor:
\begin{itemize}
    \item Los acuerdos de licenciamiento de datos (licensing deals) se volverán algo habitual, como lo demuestra la demanda por derechos de autor entre el New York Times y las empresas de IA, que está impulsando este proceso
    \item Las empresas buscan activamente adquirir o asociarse con compañías que ``poseen tesoros de datos''
    \item Los datos sintéticos (synthetic data) pueden aportar cierta mejora al entrenamiento, pero su uso excesivo provoca el \textbf{colapso de modo} (mode collapse), en el que las salidas del modelo se vuelven monótonas
    \item La investigación en modelos de lenguaje pequeños (Small Language Models) tiene un panorama prometedor: entrenar modelos pequeños para usos específicos con menos datos
\end{itemize}

\subsection{RAG: una herramienta poderosa contra las alucinaciones}

\textbf{Pregunta}: ¿Existe un conjunto de datos centralizado de alucinaciones que pueda usarse para prevenir alucinaciones futuras?

La respuesta del profesor fue reveladora: entrenar directamente con datos de alucinaciones puede resultar contraproducente, ya que el modelo podría aprender a generar alucinaciones \textbf{más verosímiles}.

\begin{importantbox}{RAG (Retrieval-Augmented Generation)}
El profesor considera que la solución más eficaz para mitigar las alucinaciones actualmente es RAG:
\begin{itemize}
    \item Mantener una base de conocimiento externa (una base de datos, un almacén vectorial, etc.) que contenga información factual confiable
    \item Enseñar al modelo a recuperar el contexto relevante de la base de conocimiento antes de generar la respuesta
    \item Este método separa de manera ingeniosa las responsabilidades:
    \begin{itemize}
        \item El LLM se encarga de lo que hace bien: generar texto fluido y coherente
        \item La base de datos se encarga de lo que hace bien: almacenar, buscar y actualizar información factual
    \end{itemize}
    \item Beneficio adicional: cuando los hechos cambian, basta con actualizar la base de datos, sin necesidad de volver a entrenar el modelo
\end{itemize}
\end{importantbox}

\subsection{¿Puede el LLM impulsar el descubrimiento científico?}

\textbf{Pregunta}: Con el crecimiento exponencial de la ventana de contexto, ¿podría el LLM impulsar avances científicos (como la cura del cáncer) al integrar resultados de investigación fragmentados?

El profesor compartió dos perspectivas:
\begin{itemize}
    \item \textbf{Vía directa}: dejar que el modelo procese directamente una gran cantidad de artículos, en busca de conexiones y patrones entre distintas áreas
    \item \textbf{Vía indirecta}: ayudar a los investigadores a digerir eficientemente enormes volúmenes de literatura (por ejemplo, resumiendo miles de artículos), acelerando el ritmo de la investigación
\end{itemize}

El profesor mencionó especialmente un caso alentador: unos investigadores aplicaron el marco de los modelos de lenguaje a la simulación del movimiento atómico (en lugar de al lenguaje natural), y el ``modelo de lenguaje atómico'' entrenado, sin haber visto nunca sal de mesa, predijo correctamente la \textbf{estructura del cristal de sal de mesa} a partir de átomos de sodio y de cloro. Esto demuestra que la capacidad predictiva de la arquitectura Transformer va mucho más allá del dominio del lenguaje.

\subsection{¿Cómo enseñar al LLM a respetar reglas?}

\textbf{Pregunta}: ¿Cómo lograr que el LLM respete reglas específicas (como las jugadas legales del ajedrez)?

El profesor propuso una solución en dos niveles:

\begin{enumerate}
    \item \textbf{A nivel del modelo}: introducir una función de penalización personalizada durante el fine-tuning, que aplique una pérdida adicional cuando el modelo genere salidas que infrinjan las reglas
    \item \textbf{A nivel del sistema}: construir fuera del modelo una arquitectura de dos capas de \textbf{capa de predicción + capa de política}. La capa de predicción (el LLM) se encarga de generar las salidas candidatas, y la capa de política (rule engine) se encarga de verificar si la salida cumple las reglas. Si no las cumple, la capa de política la rechaza y le pide al LLM que la genere de nuevo
\end{enumerate}

\begin{knowledgebox}{Arquitectura estándar de los sistemas de LLM de producción}
El profesor señala que casi todos los sistemas de aprendizaje automático de producción que ha visto —no solo los de LLM, sino también los sistemas de aprendizaje automático tradicionales anteriores a los LLM— adoptan una arquitectura de dos capas de \textbf{modelo de predicción + capa de política}. La capa de política puede bloquear, rechazar o promover (promote) determinadas salidas, con lo cual convierte un sistema esencialmente aleatorio (stochastic system) en algo controlable y predecible.
\end{knowledgebox}

\subsection{Resumen de la sección}

La sesión de preguntas y respuestas abordó varios temas de vanguardia en el ámbito de los LLM: el agotamiento de datos impulsa la investigación en datos sintéticos y la tendencia hacia modelos pequeños; RAG mitiga eficazmente las alucinaciones mediante una base de conocimiento externa; las aplicaciones científicas del LLM tienen un panorama prometedor pero requieren una evaluación cuidadosa; las restricciones de reglas en los sistemas de producción requieren un doble resguardo, tanto dentro como fuera del modelo.
